% Library CV energy_fr — transcription of writing-samples/cvs/Moumen_Nermine_CV_Energy.pdf (CV06, 2026-10-10)
\documentclass[10pt]{article}
\usepackage[letter]{nermine-cv}
\begin{document}
\cvheader{Moumen Nermine}{Étudiante en dernière année du cycle ingénieur en Chimie – Génie des procédés, énergie et environnement}{nermine.moumen@insat.ucar.tn}{+21627040414}{linkedin.com/in/nermine-moumen}

\section{FORMATION}
\entry{Institut National des Sciences Appliquées et de Technologie (INSAT)}{2022–Présent}
\summary{Étudiante en dernière année de Chimie Industrielle (spéc. Procédés et Environnement) avec une formation solide en thermodynamique, transferts thermiques, bilans matière et énergie et énergies renouvelables, complétée par la simulation de procédés et l'analyse de cycle de vie. À la recherche d'un PFE en décarbonation et efficacité énergétique de l'industrie.}

\section{EXPÉRIENCE PROFESSIONNELLE}
\entry{SEGULA Technologies - Stage en Modélisation Électrochimique des Batteries}{06/2026–08/2026}
\begin{bul}
\item Modélisation d'une cellule \textbf{Li-ion} (NMC811/Graphite-SiOx) via le Single Particle Model (SPM) implémenté sous \textbf{MATLAB}.
\item Résolution des équations aux dérivées partielles par \textbf{FVM} avec intégration temporelle de Crank-Nicolson.
\item Optimisation du maillage et des coefficients de diffusion via une analyse de données, réduisant l'erreur \textbf{RMSE de 80 \%}.
\item Validation des résultats par une approche itérative et une analyse critique avec les données expérimentales et \textbf{COMSOL}.
\end{bul}
\entry{Control Union - Stage en Contrôle Qualité}{08/2025–09/2025}
\begin{bul}
\item Analyse de lubrifiants (viscosité et métaux lourds par \textbf{ICP-OES}) avec suivi de la conformité \textbf{ISO 17025}.
\item Documentation technique et rédaction de rapports d'analyse destinés aux clients, avec rigueur et traçabilité.
\end{bul}
\entry{ETAP - Stage en traitement et caractérisation des eaux industrielles}{07/2024–08/2024}
\begin{bul}
\item Analyses physico-chimiques des eaux (pH, salinité, bicarbonates, sulfates, chlorures) par chromatographie ionique et ICP-AES.
\item Étude des méthodes de traitement des eaux pétrolières par échange d'ions et par coagulation-floculation.
\end{bul}
\section{EXPÉRIENCE COMPLÉMENTAIRE}
\entry{Cascade - Enseignante de Chimie en Ligne (Temps partiel)}{2024–Présent}
\begin{bul}
\item Accompagnement individualisé en chimie pour des élèves de Terminale française, couvrant les concepts du niveau examen.
\end{bul}

\section{PROJETS ACADÉMIQUES}
\project{Analyse techno-environnementale de la production d'éthylène : vapocraquage vs déshydratation du bioéthanol}
\begin{bul}
\item Simulation comparative de deux voies de production sous \textbf{Aspen Plus V14}, avec établissement des \textbf{bilans matière et énergie} et des diagrammes de flux (PFD) du procédé, pour évaluer la faisabilité technique de chaque filière.
\item \textbf{Analyse de Cycle de Vie} cradle-to-gate sous \textbf{OpenLCA} pour comparer les impacts environnementaux des deux voies.
\end{bul}
\project{Simulation du procédé de production d'oxyde d'éthylène sous Aspen HYSYS V14}
\begin{bul}
\item Simulation du procédé en utilisant la thermodynamique NRTL-RK et la cinétique de Langmuir-Hinshelwood, avec établissement des bilans de matière et d'énergie et \textbf{dimensionnement des équipements clés}.
\item Optimisation des conditions opératoires par analyse de sensibilité, atteignant \textbf{96 \% de récupération} et \textbf{97,14 \% de pureté}.
\end{bul}
\project{Modélisation de Piles à Combustible PEM et Électrolyseurs PEM sous COMSOL et ANSYS Fluent}
\begin{bul}
\item Modélisation de piles à combustible : écoulement diphasique, effets thermiques, engorgement de la couche de diffusion.
\item Modélisation d'un \textbf{électrolyseur PEM 3D} : conservation de charge, identification de la réaction limitante (OER).
\end{bul}

\section{COMPÉTENCES TECHNIQUES}
\begin{bul}
\item \textbf{Simulation :} Aspen Plus et HYSYS, COMSOL Multiphysics\sep\textbf{Programmation :} \textit{MATLAB, Mathematica, Python}
\item \textbf{Données \& Bureautique :} Pack Microsoft Office, OpenLCA\sep\textbf{Normes :} ISO 9001, ISO 14001, ISO 45001, ISO 17025, ISO 14040
\item \textbf{Techniques de Laboratoire :} Chromatographie (HPLC, GC, CCM), Spectrométrie (MS, ICP-OES), Spectroscopie (RMN, FTIR, UV-Vis), Diffraction (XRD), Microscopie (MEB, MET)
\end{bul}

\section{ACTIVITÉS EXTRASCOLAIRES}
\textbf{IEEE INSAT} — Chef d'organisation de TRSYP 3.0 ; Responsable des activités sociales – MENASYP ; Présidente de l'EMBS.\\
\textbf{THEATRO INSAT} — Membre de l'équipe de danse.\par

\section{LANGUES ET CERTIFICATS}
\twocol{Langues}{Anglais (Avancé), Français (Bilingue), Arabe (Natif).}{Certificats}{Aspen HYSYS – MATLAB}
\end{document}
